% ==========================================
% LatexRefinement Step Output: step4-2026-09-23-wednesday-week2-algebra-1-offset-final.tex
% Model: gemini-3.6-flash
% Temperature: 0.35
% TopP: 0.9
% TopK: 10
% MaxOutputTokens: 65535
% ThinkingBudget: 4096
% ThinkingLevel: HIGH
% Prompt Tokens: 52’384
% Candidates Tokens: 20’975 (inkl. Thinking Tokens)
% Cached Tokens: 0
% Processed on: 2026-10-03 21:34:13
% ==========================================

% Primary Language: English
% End of the video: 00:42:33

\lecturechapter[2]{Wednesday}{23 Sep}{23 September 2026}{Category Theory Basics, Integral Domains, and Groups of Units}
% Old chapter title: Category Theory Basics \& Integral Domains

\begin{overview}[Lecture Overview \& Context]
This lecture concludes the foundational discussion of category theory before returning to core ring theory and modular arithmetic:
\begin{itemize}
    \item \textbf{Category Theory Basics:} Abstract definitions of objects, morphism sets (hom-sets), composition laws, identity morphisms, isomorphisms, monomorphisms (left-cancellation / injectivity), and epimorphisms (right-cancellation / surjectivity) along with categorical duality.
    \item \textbf{Integral Domains:} Zero divisors, formal definition of an integral domain (including the mandatory exclusion of the zero ring), and the proof that $\mathbb{Z}/m\mathbb{Z}$ is an integral domain if and only if $m$ is prime.
    \item \textbf{Divisibility and Units:} Divisibility in arbitrary commutative rings, units, and precise linguistic distinctions (\qt{the units}, \qt{the unit}, \qt{a unit}).
    \item \textbf{Groups of Units \& Modular Isomorphisms:} Proof that the set of units $U(R)$ forms a multiplicative group, characterization of $U(\mathbb{Z}/m\mathbb{Z})$ via coprimality, and explicit proof that $U(\mathbb{Z}/7\mathbb{Z}) \cong (\mathbb{Z}/6\mathbb{Z}, +)$ as a cyclic group generated by $\bar{3}$.
\end{itemize}
\end{overview}

\section{Administrative Remarks}

\begin{speech}[00:00:00 - 00:01:26]
Okay, so we start. So, I wanted to make this comment: there were several questions about the course. The first one is if you want to watch the lectures, you have to be enrolled. So, if you're not enrolled, it will stop you. And there's not much I can do about that, or if there is, I don't know what to do, but the idea is you should enroll in the course if you want to watch the lectures.

Other people had some difficulty uploading the problem sets. Maybe you have to be enrolled in the course for that, too --- I don't know ---, but many people were successful, so if you have difficulty, find someone who's successful, because I don't know how to solve the problem in another way.

Okay, a question that's more about the actual course was that every week there are these suggested pages to read in the book, and the question is: how do they relate to the course? That's a good question.

And the answer is that, roughly speaking, I do that so you can find the parallel material on which I'm lecturing in the book. And I don't always follow exactly the order there, I don't follow exactly the notation, and I don't follow exactly the text of the definitions. But I do something that's logically equivalent, so it can be helpful. As far as matters like the exam, I will only test you on material that has appeared in the lecture.

Okay? So, this is my best statement about this so far. Okay.
\end{speech}

\section{Category Theory: Abstract Definitions}

\subsection{Definition and Axioms of a Category}

\begin{speech}[00:01:26 - 00:02:50]
So, the last topic yesterday was categories, and I wanted to say a few things.

When I had this discussion about categories yesterday, it was mixed in with all the examples, so I want to just very quickly give the pure abstract definition and two small points. But this course is not about category theory. It's a topic here because it's natural to introduce when you talk about these algebraic objects, so, anyway, the course is not about category theory.

So, a category has two types of data: it has objects, as we said, and it has homs. So, $A$ can be an object, and for homs, if you pick any two objects, you get a hom space, sometimes called morphisms.

And the abstract properties of this hom space --- I said them, but I will repeat ---: first, that it's a set, that's a technical property. And that's to avoid some contradictions in foundations, which I'm not going to discuss now.

The second is that there's a composition law.

And the third is that there are identity elements --- existence of identity.

So... \inlinemetanote{writes on board} Okay, that's the data. So, is there any question about this data?
\end{speech}

\begin{content}[Constituents of a Category: Data]
\recap{Categories were introduced the day before, mixed in with examples (\cref{def:category-foundations}, \cref{def:category-full}). The lecturer now repeats the abstract definition on its own.}

\begin{definition}[Category Data]\label[definition]{def:category-data-w2wed}
A \newterm{category} $\mathcal{C}$ consists of the following mathematical data:
\begin{itemize}
    \item \textbf{Objects:} A class of objects, denoted $\operatorname{Obj}(\mathcal{C})$. We write $A \in \operatorname{Obj}(\mathcal{C})$ (or simply $A \in \mathcal{C}$) to denote that $A$ is an object of $\mathcal{C}$.
    \item \textbf{Morphism Sets (Homs):} For any pair of objects $A, B \in \operatorname{Obj}(\mathcal{C})$, a set of morphisms from $A$ to $B$, denoted:
    \[
    \operatorname{Hom}(A, B) \quad \text{or} \quad \operatorname{Hom}_{\mathcal{C}}(A, B).
    \]
    An element $f \in \operatorname{Hom}(A, B)$ is represented diagrammatically as an arrow $f \colon A \to B$.
    \item \textbf{Composition Law:} For any three objects $A, B, C \in \operatorname{Obj}(\mathcal{C})$, a binary composition mapping:
    \begin{align*}
        \circ \colon \operatorname{Hom}(B, C) \times \operatorname{Hom}(A, B) &\to \operatorname{Hom}(A, C), \\
        (g, f) &\mapsto g \circ f.
    \end{align*}
    \item \textbf{Identities:} For each object $A \in \operatorname{Obj}(\mathcal{C})$, a designated morphism called the identity morphism on $A$:
    \[
    \operatorname{Id}_A \in \operatorname{Hom}(A, A) \quad (\text{or } 1_A \in \operatorname{Hom}(A, A)).
    \]
\end{itemize}
\end{definition}

\begin{steps}[Foundational Remark on Hom-Sets]
The requirement that $\operatorname{Hom}(A, B)$ forms a legitimate set (rather than a proper class) is a standard set-theoretic convention defining a \emph{locally small category}. This prevents Russell-style self-referential paradoxes when indexing collections of arrows.
\end{steps}
\end{content}

\begin{speech}[00:02:50 - 00:03:45]
And then there's the axioms. The first one is that the composition law is associative, which means, you know, if you have three things, you can compose them in the two different ways. I hope I don't have to say more than that.

And then the second is the axioms about the identity, which is roughly: if I do this with $f$, this is $f$, and the other way around also. The identity doesn't change $f$. So, we discussed this last time.
\end{speech}

\begin{content}[Associativity and Identity Axioms]
\begin{definition}[Category Axioms]\label[definition]{ax:category-axioms-w2wed}
The data of a category $\mathcal{C}$ must satisfy the following two fundamental axioms:
\begin{enumerate}[label=\textbf{(\arabic*)}]
    \item \textbf{Associativity of Composition:} For all objects $A, B, C, D \in \operatorname{Obj}(\mathcal{C})$ and all morphisms $f \in \operatorname{Hom}(A, B)$, $g \in \operatorname{Hom}(B, C)$, and $h \in \operatorname{Hom}(C, D)$, composition is associative:
    \begin{equation}\label{eq:cat-assoc}
    (h \circ g) \circ f = h \circ (g \circ f).
    \end{equation}
    \item \textbf{Identity Axioms (Neutrality):} For any morphism $f \in \operatorname{Hom}(A, B)$, composition with the respective identity morphisms leaves $f$ unchanged:
    \begin{align}
        f \circ \operatorname{Id}_A &= f, \label{eq:cat-id-right} \\
        \operatorname{Id}_B \circ f &= f. \label{eq:cat-id-left}
    \end{align}
\end{enumerate}
\end{definition}

{\centering
\begin{tikzpicture}[scale=1.4, >=Stealth, auto]
    \node (A) at (0, 0) {$A$};
    \node (B) at (2, 1.3) {$B$};
    \node (C) at (4, 0) {$C$};

    \draw[->, thick] (A) to node {$f$} (B);
    \draw[->, thick] (B) to node {$g$} (C);
    \draw[->, thick, MidnightBlue] (A) to node[below] {$g \circ f$} (C);

    \path[->, thick] (A) edge[out=150, in=210, looseness=4] node[left] {$\operatorname{Id}_A$} (A);
    \path[->, thick] (B) edge[out=60, in=120, looseness=4] node[above] {$\operatorname{Id}_B$} (B);
    \path[->, thick] (C) edge[out=330, in=30, looseness=4] node[right] {$\operatorname{Id}_C$} (C);
\end{tikzpicture}\par}
\end{content}

\subsection{Categorical Morphisms: Isomorphisms, Monomorphisms, and Epimorphisms}

\begin{speech}[00:03:45 - 00:05:35]
So, that's the abstract definition, so there's really not that much here. And we discussed this last time, and we had lots of examples.

And so you could say: well, why do we want to do this? And as I said, there's a setting where many concepts that we talk about when we talk about groups, rings, fields, vector spaces --- they're categorical concepts, meaning that they only depend on this structure. So, it simplifies thinking to consider them all together.

So, I already told you about isomorphism. The categorical isomorphism is that if I have this arrow from $A$ to $B$ --- this picture means this: $f \in \operatorname{Hom}(A, B)$, okay? That's why in category theory, you can say this, or you can say that, and you mean the same thing --- that this $f$ is an isomorphism --- that's a definition --- if and only if there exists $g$ the other way, $g \colon B \to A$, such that the two compositions are identities.

And if you want to say which identity, you have to be careful. So, $g$ starts with $B$, so if you do $g$ first, then $f$ back, this should be the identity in $B$, and the other is the identity in $A$.

So, that's the abstract definition of isomorphism. And in all of the categories we discussed, that coincides with the notion of isomorphism in the category that you already know, like, for example, isomorphism of vector spaces.
\end{speech}

\begin{content}[Categorical Isomorphism]
\recap{The categorical isomorphism was defined the day before (\cref{def:categorical-isomorphism}); the lecturer restates it.}

\begin{definition}[Isomorphism]\label[definition]{def:categorical-isomorphism-w2wed}
Let $\mathcal{C}$ be a category and let $f \in \operatorname{Hom}(A, B)$ (written $f \colon A \to B$). The morphism $f$ is an \newterm{isomorphism} if and only if there exists a morphism $g \in \operatorname{Hom}(B, A)$ (written $g \colon B \to A$) such that:
\begin{align}
    f \circ g &= \operatorname{Id}_B, \label{eq:isom-right-inverse} \\
    g \circ f &= \operatorname{Id}_A. \label{eq:isom-left-inverse}
\end{align}
If such a morphism $g$ exists, it is uniquely determined, called the \newterm{inverse} of $f$, and denoted $f^{-1} := g$. Two objects $A$ and $B$ are said to be \newterm{isomorphic} (written $A \cong B$) if there exists an isomorphism between them.
\end{definition}

{\centering
\begin{tikzpicture}[scale=1.5, >=Stealth, auto]
    \node (A) at (0,0) {$A$};
    \node (B) at (3,0) {$B$};

    \draw[->, thick, RoyalBlue] (A) to[bend left=25] node {$f$} (B);
    \draw[->, thick, BurntOrange] (B) to[bend left=25] node {$g = f^{-1}$} (A);

    % caption below the label of the lower arc (arc dips to y = -0.33, its label to about y = -0.55)
    \node at (1.5, -1.15) {\footnotesize $f \circ g = \operatorname{Id}_B \quad \text{and} \quad g \circ f = \operatorname{Id}_A$};
\end{tikzpicture}\par}

\begin{steps}[Consistency across Concrete Categories]
In concrete categories where objects are sets endowed with structure and morphisms are structure-preserving maps (homomorphisms):
\begin{itemize}
    \item In $\mathbf{Set}$, isomorphisms are precisely bijective functions.
    \item In $\mathbf{Grp}$, $\mathbf{Ring}$, and $\mathbf{Vect}_K$, isomorphisms are bijective homomorphisms (for groups: \cref{prop:isomorphism-inverse-characterization}).
\end{itemize}
The power of the categorical definition is that it never inspects elements $x \in A$; it characterizes invertibility purely through diagrammatic composition.
\end{steps}
\end{content}

\begin{hidden}
% timestamp: 00:05:35
% topic: Abstract Definition of Category, Axioms, and Isomorphism
% board_state: def:category-data-w2wed, ax:category-axioms-w2wed, def:categorical-isomorphism-w2wed
% next_goal: Introduce monomorphisms and epimorphisms as categorical analogues of injectivity and surjectivity
% open_loops: none
\end{hidden}

\begin{speech}[00:05:35 - 00:07:20]
There are two other notions, and that's the last thing I want to say. And this is the notion of monomorphism.

So, $f$ from $A$ to $B$ is a monomorphism --- that's a long word, that's a definition --- if for all $g_1$ arrows from another object, $X$, to $A$, and another one, $g_2$ --- so we take any two ---, then the composition $f \circ g_1 = f \circ g_2$ implies $g_1 = g_2$. I hope I wrote that clearly here, okay?

So, it's a little bit complicated to say. If I have an arrow in the category, it may or may not be a monomorphism. I'm defining a monomorphism. And the condition of being a monomorphism means that whenever I take an object $X$ and map it two ways to $A$, the composition with $f$ being equal implies they were equal to start with. Of course, if they're equal to start with, the composition is equal, but it's the other implication. That's the definition of a monomorphism.

And, you know, you have to think about this and unwind it. But once you have this definition, you can imagine: what is a monomorphism in all of those categories you know? And what do you think it is?
\end{speech}

\begin{speech}[00:07:20 - 00:08:15]
So, I leave this for you to think about, but roughly speaking, monomorphism is being injective.

Just imagine they're sets. If this is injective, and you take any two maps, you can compare them after you do $f$. So, you have to think about it, but this is the definition.

And once you take this definition, you can say: okay, I can take the category of sets and functions, and what is a monomorphism there? And the answer is: injective. That's the \qt{mono} in monomorphism.

So, this is a categorical formulation of being injective.

I don't know if you find this appealing or not, but that's what this discussion is about. Because the nice thing about this formulation is that it applies to all examples of categories. Every single one we discussed, and all the other ones you can think about, there's a notion of monomorphism.
\end{speech}

\begin{content}[Left-Cancellable Morphisms: Monomorphisms]
\begin{definition}[Monomorphism]\label[definition]{def:monomorphism}
Let $\mathcal{C}$ be a category. A morphism $f \in \operatorname{Hom}(A, B)$ (or $f \colon A \to B$) is called a \newterm{monomorphism} (often abbreviated as a \emph{monic morphism} or \emph{mono}) if it is left-cancellable: for every object $X \in \operatorname{Obj}(\mathcal{C})$ and all pairs of parallel morphisms $g_1, g_2 \in \operatorname{Hom}(X, A)$,
\begin{equation}\label{eq:monomorphism-cond}
f \circ g_1 = f \circ g_2 \implies g_1 = g_2.
\end{equation}
\end{definition}

{\centering
\begin{tikzpicture}[scale=1.5, >=Stealth, auto]
    \node (X) at (0,0) {$X$};
    \node (A) at (2.2,0) {$A$};
    \node (B) at (4.4,0) {$B$};

    \draw[->, thick] (X) to[bend left=25] node {$g_1$} (A);
    \draw[->, thick] (X) to[bend right=25] node[below] {$g_2$} (A);
    \draw[->, thick, ForestGreen] (A) to node {$f$} (B);

    % caption below the g_2 label (lower arc dips to y = -0.24, its label to about y = -0.5)
    \node at (2.2, -1.0) {\footnotesize $f \circ g_1 = f \circ g_2 \implies g_1 = g_2 \quad (\text{Categorical Injective})$};
\end{tikzpicture}\par}

\begin{steps}[Why Monomorphisms Characterize Injectivity]
In the category of sets ($\mathbf{Set}$), suppose $f \colon A \to B$ is a monomorphism. To test whether $f$ is injective on points $a_1, a_2 \in A$, let $X = \{*\}$ be a one-element set and define two maps $g_1, g_2 \colon \{*\} \to A$ by $g_1(*) = a_1$ and $g_2(*) = a_2$.
Then $(f \circ g_1)(*) = f(a_1)$ and $(f \circ g_2)(*) = f(a_2)$. Thus:
\[
f(a_1) = f(a_2) \implies f \circ g_1 = f \circ g_2 \implies g_1 = g_2 \implies a_1 = a_2.
\]
Hence, left-cancellation with respect to arbitrary test objects $X$ precisely recovers element-wise injectivity without mentioning elements.
\end{steps}
\end{content}

\begin{speech}[00:08:15 - 00:09:20]
And this one has a sibling, or a partner --- I don't know their relationship ---, which is... whenever you talk about injective, you can also talk about surjective, and the categorical term there is epimorphism.

So, you can define epimorphism. So, if you have $f \colon A \to B$, it is an epimorphism if --- it's a definition --- for all...

And this is now about this kind of feeling of mathematics: how do you think I'm going to define the opposite notion of injective? Like here we had injective by $X$ mapping to $A$. So, what do you think I should do here?

You should guess. Guessing is like one of the most important tools of a mathematician.

Yeah?
\end{speech}

\begin{student}[Student Answer]
Well, we do the same, except we also switch like we compose\dots
\end{student}

\begin{speech}[00:09:25 - 00:10:35]
This is a long answer. Here we had $X$ mapping to $A$, so the opposite of that is mapping from $B$. It's like the other side.

So, for every $B \xrightarrow[g_2]{g_1} X$, then... So, this is the opposite of that, by just flipping everything to the other side. So, for every one of these, then we can consider $f$...

In that one it was the $g$'s first, then $f$; now it's $f$ first, then the $g$'s: $g_1 \circ f = g_2 \circ f \implies g_1 = g_2$.

So, that's the categorical notion of epimorphism, and it's somehow dual to monomorphism. If you like this world, it almost makes epimorphism and monomorphism the same definition up to this flipping. And, you know, thinking about things categorically is good for such understanding.

If you look at this definition, and you look at it for the examples we know, this is surjectivity.

Okay, so that's the end of this discussion about categories. If you like it, you can go read more. As I said, this course is not a course on categories. There's a bunch of other things that I would say here, but I'm not going to, because, yeah, maybe I will at the end if I have time.
\end{speech}

\begin{content}[Right-Cancellable Morphisms: Epimorphisms]
\begin{definition}[Epimorphism]\label[definition]{def:epimorphism}
Let $\mathcal{C}$ be a category. A morphism $f \in \operatorname{Hom}(A, B)$ (or $f \colon A \to B$) is called an \newterm{epimorphism} (often abbreviated as an \emph{epic morphism} or \emph{epi}) if it is right-cancellable: for every object $X \in \operatorname{Obj}(\mathcal{C})$ and all pairs of parallel morphisms $g_1, g_2 \in \operatorname{Hom}(B, X)$,
\begin{equation}\label{eq:epimorphism-cond}
g_1 \circ f = g_2 \circ f \implies g_1 = g_2.
\end{equation}
\end{definition}

{\centering
\begin{tikzpicture}[scale=1.5, >=Stealth, auto]
    \node (A) at (0,0) {$A$};
    \node (B) at (2.2,0) {$B$};
    \node (X) at (4.4,0) {$X$};

    \draw[->, thick, BrickRed] (A) to node {$f$} (B);
    \draw[->, thick] (B) to[bend left=25] node {$g_1$} (X);
    \draw[->, thick] (B) to[bend right=25] node[below] {$g_2$} (X);

    % caption below the g_2 label (lower arc dips to y = -0.24, its label to about y = -0.5)
    \node at (2.2, -1.0) {\footnotesize $g_1 \circ f = g_2 \circ f \implies g_1 = g_2 \quad (\text{Categorical Surjective})$};
\end{tikzpicture}\par}
\end{content}

\begin{aiNote}[Epimorphisms Need Not Be Surjective]
The lecturer says that for the examples we know, epimorphism means surjectivity. This holds in $\mathbf{Set}$, $\mathbf{Grp}$ and $\mathbf{Vect}_K$, but not in $\mathbf{Ring}$: the inclusion $\iota \colon \mathbb{Z} \hookrightarrow \mathbb{Q}$ is an epimorphism without being surjective. Indeed, if $g_1, g_2 \colon \mathbb{Q} \to R$ are ring homomorphisms with $g_1 \circ \iota = g_2 \circ \iota$, then for $b \neq 0$ the element $g_k(b)$ is invertible with inverse $g_k(\frac{1}{b})$, so $g_1(\frac{a}{b}) = g_1(a) \, g_1(b)^{-1} = g_2(a) \, g_2(b)^{-1} = g_2(\frac{a}{b})$. For monomorphisms the slogan is safe: in $\mathbf{Set}$, $\mathbf{Grp}$, $\mathbf{Ring}$ and $\mathbf{Vect}_K$ they are exactly the injective homomorphisms.
\end{aiNote}

\begin{insight}[Categorical Duality: Reversing Arrows]
Monomorphisms and epimorphisms represent our first concrete encounter with \newterm{categorical duality}. If we construct the opposite category $\mathcal{C}^{\text{op}}$ by reversing all morphism arrows ($f^{\text{op}} \colon B \to A$ for every $f \colon A \to B$), then:
\[
f \text{ is a monomorphism in } \mathcal{C} \iff f^{\text{op}} \text{ is an epimorphism in } \mathcal{C}^{\text{op}}.
\]
Every theorem or statement about monomorphisms immediately yields an identical, dual statement about epimorphisms simply by reversing the order of all composition arrows.
\end{insight}

\section{Ring Theory: Zero Divisors and Integral Domains}

\subsection{Zero Divisors and the Definition of an Integral Domain}

\begin{speech}[00:10:35 - 00:12:35]
So, let us return to rings now. This was kind of a foundational side note.

So, we have this ring $(R, +, \cdot, 0, 1)$.

And the rings here, there's kind of a different flavor of rings: there are the rings that are like $\mathbb{Z}$, where when you have things that are non-zero and you multiply them, they stay non-zero. That's one of the principles of $\mathbb{Z}$. But it's not true for rings like $\mathbb{Z}/6\mathbb{Z}$, this ring.

I can take two elements here, like the residue class of $2$ times the residue class of $3$, neither of which is zero. So, this is not zero, and that is not zero. So, I take two elements that are legitimately not zero, but when I multiply these residue classes, this is $6$, and that is $0$.

So, that's a different type of multiplication. It's a multiplication where when you multiply non-zero things, you get zero, which is unexpected from the point of view of the integers.

And this is also related to the fact that there's no hope of having a multiplicative inverse if you're mapping things to zero... Okay, so there are these two differences, and there's a definition for this:
\end{speech}

\begin{content}[Zero Divisors in Modular Arithmetic]
Let $(R, +, \cdot, 0, 1)$ be a commutative ring with unity $1 \neq 0$.

\begin{example}[Multiplication in \texorpdfstring{$\mathbb{Z}/6\mathbb{Z}$}{Z/6Z}]\label[example]{ex:z6z-zero-divisors}
In the ring of residue classes modulo $6$, denoted $\mathbb{Z}/6\mathbb{Z}$ (or $\mathbb{Z}_6$):
\begin{align*}
    \bar{2} &\neq \bar{0}, \\
    \bar{3} &\neq \bar{0}.
\end{align*}
However, their product in the ring is:
\[
\bar{2} \cdot \bar{3} = \overline{2 \cdot 3} = \bar{6} = \bar{0} \in \mathbb{Z}/6\mathbb{Z}.
\]
\end{example}

\begin{steps}[Zero Divisors vs. Invertibility]
Non-zero elements whose product is zero are called \newterm{zero divisors}. If an element $a \in R$ is a zero divisor (with $a \cdot b = 0$ for some $b \neq 0$), it can never possess a multiplicative inverse $a^{-1}$:
\[
b = 1 \cdot b = (a^{-1} \cdot a) \cdot b = a^{-1} \cdot (a \cdot b) = a^{-1} \cdot 0 = 0,
\]
which contradicts $b \neq 0$. Hence, rings containing zero divisors fail to exhibit regular arithmetic cancellation.
\end{steps}
\end{content}

\begin{speech}[00:12:35 - 00:14:05]
A ring $R$ is an integral domain... So, it's a definition of this term \qt{integral domain}. And some authors would just say it's a domain, and also if you speak about rings, very often mathematicians will say it's a domain... if whenever I have $a$ and $b$, neither of which is zero --- so maybe I'll just be really explicit --- whenever I have $a \neq 0$ and $b \neq 0$, this implies the product is not equal to zero ($a \cdot b \neq 0$).

So, that's the definition. It's a really basic definition in ring theory, and it separates this type of ring from that type of ring. Is that clear?

Fields, for example, are always domains, because you have the multiplicative inverse, if you think about it, okay?

So, we can ask...
\end{speech}

\begin{content}[Definition of an Integral Domain (Initial Board State)]
\begin{definition}[Integral Domain / Domain]\label[definition]{def:integral-domain-prelim}
A commutative ring $R$ is called an \newterm{integral domain} (or simply a \newterm{domain}) if for all $a, b \in R$:
\begin{equation}\label{eq:integral-domain-cond}
a \neq 0 \quad \text{and} \quad b \neq 0 \implies a \cdot b \neq 0.
\end{equation}
Equivalently, taking the contrapositive:
\[
a \cdot b = 0 \implies a = 0 \quad \text{or} \quad b = 0.
\]
\end{definition}

\begin{example}[Fields are Domains]\label[example]{ex:fields-are-domains}
Every field $F$ is an integral domain: if $a \neq 0$ and $a \cdot b = 0$, then multiplying by the inverse $a^{-1} \in F$ immediately yields:
\[
b = 1 \cdot b = (a^{-1} a) b = a^{-1} (a b) = a^{-1} \cdot 0 = 0.
\]
\end{example}
\end{content}

\begin{speech}[00:14:05 - 00:14:42]
So, see, I did that already there: I called it a domain, not an integral domain, because it takes a lot of syllables to say \qt{integral domain}. Formally, this is defined as an integral domain; colloquially, almost everyone says a domain.

So, these terms, monomorphism and epimorphism, have you ever heard these terms before? In which context?
\end{speech}

\begin{student}[Student Question/Answer]
Linear algebra.
\end{student}

\begin{speech}[00:14:42 - 00:14:52]
Ah, okay. And they were just used synonymously with injection and surjection?
\end{speech}

\begin{student}[Student Answer]
We were introduced to them, but then we went back to using injection and surjection.
\end{student}

\begin{speech}[00:14:52 - 00:14:58]
Okay, yeah.
\end{speech}

\begin{student}[Student Answer]
So, they were like injective and surjective homomorphisms.
\end{student}

\begin{speech}[00:14:58 - 00:15:20]
Yes, exactly, yes. Injective and surjective homomorphisms. That's precisely what this categorical thing says, right? Because in the category of groups, the arrows are homomorphisms, and when a homomorphism is a monomorphism, it's an injective homomorphism.

In my experience with undergraduates, some really like the idea of categories and really like to think in that language, and others like the examples, so either way is okay.
\end{speech}

\subsection{Characterization of \texorpdfstring{$\mathbb{Z}/m\mathbb{Z}$}{Z/mZ} as an Integral Domain}

\begin{speech}[00:15:20 - 00:16:24]
All right, so here we come to the first proposition, which is something my guess is you are already familiar with, but we do it again. So, we take this ring $\mathbb{Z}/m\mathbb{Z}$, and the claim is that this $\mathbb{Z}/m\mathbb{Z}$ is a domain if and only if --- and what am I going to write here? Yeah?
\end{speech}

\begin{student}[Student Answer]
$m$ is prime.
\end{student}

\begin{speech}[00:16:04 - 00:16:24]
$m$ is prime. $m$ is prime.

Oh, yeah, I should say one thing about domains which is important, I didn't say, and it's kind of annoying: part of the definition... I actually should put that in.
\end{speech}

\begin{speech}[00:16:24 - 00:17:15]
It's a super annoying thing to talk about, but I talk about any ring here, so my definition has to deal with the zero ring, the one that, you know, is unpleasant to think about. And you could think about: what would this definition be for the zero ring? Because for all $a$ and $b$, there are no $a$'s and $b$'s that are not equal to zero. So, formally this statement in the formal deductive logical sense is true for the zero ring.

So, it is part of the definition: a ring $R$ which is not the zero ring ($R \neq \{0\}$) is a domain. Zero rings are not allowed to be domains, by fiat, in the definition.
\end{speech}

\begin{content}[Correction: Exclusion of the Zero Ring]
The lecturer corrects \cref{def:integral-domain-prelim} by adding condition (i):

\begin{definition}[Integral Domain (Formal Standard Definition)]\label[definition]{def:integral-domain}
A commutative ring $R$ is an \newterm{integral domain} (or \newterm{domain}) if:
\begin{enumerate}[label=\textbf{(\roman*)}]
    \item $R \not\cong \{0\}$ (that is, $1_R \neq 0_R$, excluding the zero ring; cf.\ \cref{prop:zero-ring}).
    \item For all $a, b \in R$:
    \[
    a \neq 0 \quad \text{and} \quad b \neq 0 \implies a \cdot b \neq 0.
    \]
\end{enumerate}
\end{definition}

\begin{steps}[Why the Zero Ring Satisfies the Conditional Vacuously]
In the zero ring $R = \{0\}$, the only element is $0$. Therefore, there exist no non-zero elements $a, b \in R$. The premise $(a \neq 0 \land b \neq 0)$ is vacuously false for all elements, making the logical implication $(a \neq 0 \land b \neq 0 \implies ab \neq 0)$ vacuously true! To avoid having the degenerate zero ring classified as a domain, mathematicians add $R \neq 0$ explicitly by fiat.
\end{steps}
\end{content}

\begin{content}[Proposition: When Residue Class Rings are Domains]
\begin{proposition}\label[proposition]{prop:z-mod-m-domain}
Let $m \in \mathbb{Z}_{>0}$ be a positive integer. The residue class ring $(\mathbb{Z}/m\mathbb{Z}, +, \cdot, \bar{0}, \bar{1})$ is an integral domain if and only if $m$ is a prime number:
\begin{equation}\label{eq:z-mod-m-equiv}
\mathbb{Z}/m\mathbb{Z} \text{ is an integral domain} \iff m \text{ is prime}.
\end{equation}
\end{proposition}
\end{content}

\begin{hidden}
% timestamp: 00:17:15
% topic: Proof of Proposition: Z/mZ is an integral domain iff m is prime
% board_state: def:integral-domain, prop:z-mod-m-domain
% next_goal: Direction 1 (m not prime implies not domain) and Direction 2 (m prime implies domain via Euclid's lemma)
% open_loops: none
\end{hidden}

\begin{proofWrapper}[Proof of Proposition \ref{prop:z-mod-m-domain}]
\begin{speech}[00:17:15 - 00:19:24]
Okay, $\mathbb{Z}/m\mathbb{Z}$ is a domain if and only if $m$ is prime. Okay, so let's try to prove this. So, proof.

This is like one of the issues I said: this issue about the zero ring is a little bit problematic everywhere in the statements, because for any statement you write, \qt{let $R$ be a ring}, you have to think --- if you want that statement to be literally true in the mathematical sense from the definitions in this class ---, you have to think about what that statement could mean for the zero ring and whether it's true. And as I said, I don't view it as a deep or really important thing to do, but in terms of bookkeeping and saying things correctly, you have to say it...

Okay, so let's try to prove this: $\mathbb{Z}/m\mathbb{Z}$ is a domain if and only if $m$ is prime.

So, \textbf{(1)} suppose $m$ is not prime.

Here we have $m$ as some positive integer ($m \in \mathbb{Z}_{>0}$). Suppose $m$ is not prime. Then this means that $m$ admits a factorization into $a \cdot b$, where $a$ is strictly greater than $1$ and strictly less than $m$, and $b$ is strictly like this ($1 < b < m$). That's what it means. If it's not prime, I can factor it, and if I can factor it, I break it into two smaller pieces.

And then if I look at this equation in residue classes, this equals that, which is equal to zero in $\mathbb{Z}/m\mathbb{Z}$; and this shows that that's not zero ($\bar{a} \neq \bar{0}$), and this shows that that's not zero ($\bar{b} \neq \bar{0}$).

Okay? So, this says that if $m$ is not prime, there is no way for it to be a domain, because if it's not prime, it factors, and those factors precisely give the violation of that property. All right?
\end{speech}

\begin{content}[Proof of Proposition: Direction 1 (Non-Prime Case)]
\textbf{Direction 1: \texorpdfstring{$m$}{m} Not Prime \texorpdfstring{$\implies \mathbb{Z}/m\mathbb{Z}$}{=> Z/mZ} Not a Domain.}

Suppose $m$ is not prime. If $m = 1$, then $\mathbb{Z}/1\mathbb{Z} = \{\bar{0}\}$ is the zero ring, which is not a domain by condition (i) of \cref{def:integral-domain} (see also \cref{rem:boundary-cases-mod-rings}). Otherwise $m > 1$ is composite and admits a proper non-trivial factorization in $\mathbb{Z}$:
\begin{align}
    m &= a \cdot b, \label{eq:composite-m-factor} \\
    1 &< a < m, \nonumber \\
    1 &< b < m. \nonumber
\end{align}
Passing to the residue classes in the quotient ring $\mathbb{Z}/m\mathbb{Z}$:
\begin{equation}\label{eq:prod-zero-divisors}
\bar{a} \cdot \bar{b} = \overline{a \cdot b} = \bar{m} = \bar{0} \in \mathbb{Z}/m\mathbb{Z}.
\end{equation}
However, because $1 < a < m$ and $1 < b < m$, neither $a$ nor $b$ is divisible by $m$. Therefore:
\begin{equation}\label{eq:nonzero-factors}
\bar{a} \neq \bar{0} \quad \text{and} \quad \bar{b} \neq \bar{0}.
\end{equation}
The existence of non-zero elements $\bar{a}, \bar{b} \in \mathbb{Z}/m\mathbb{Z}$ satisfying $\bar{a} \cdot \bar{b} = \bar{0}$ demonstrates that $\mathbb{Z}/m\mathbb{Z}$ possesses zero divisors, so $\mathbb{Z}/m\mathbb{Z}$ is not an integral domain.
\end{content}

\begin{speech}[00:19:24 - 00:21:50]
So, then the other thing is suppose it is prime. Suppose $m = p$ is prime.

Then, we let $a$ and $b$ be integers such that $\bar{a} \cdot \bar{b} = \bar{0}$ in $\mathbb{Z}/p\mathbb{Z}$. All right?

What are we going to do here? This means that if I translate this equation on residue classes to integers, this means that $a \cdot b$ is equal to some integer $d \cdot p$.

So, I am using this notation and arithmetic in $\mathbb{Z}/m\mathbb{Z}$ freely here, because I assume you've seen this before. Is this true? If it's not, you should tell me. All right. If you haven't seen it, you should study it.

So, if I take these two residue classes and their product is zero, that exactly means that the product is divisible by $p$. Is that clear?

And then this means that $p$ divides $a \cdot b$, because $a \cdot b = d \cdot p$.

And then there's the principle of Euclid we discussed in the first week that says either $p$ divides $a$ or $p$ divides $b$. It wasn't such a long time ago, it was just last week.

And this means if $p \mid a$, then that means $\bar{a} = \bar{0}$; and if $p \mid b$, that means $\bar{b} = \bar{0}$.

So, we've proven that if the product is zero in $\mathbb{Z}/p\mathbb{Z}$, the only way that happens is one of them is zero, or the other is zero. And that's logically equivalent to this thing: if they were both not zero, then the product is not zero. Or the logical equivalent, which is the contrapositive, is that if the product is zero, then one of the factors must be zero, which we've established. So, this proposition is proven.
\end{speech}

\begin{content}[Proof of Proposition: Direction 2 (Prime Case via Euclid's Lemma)]
\textbf{Direction 2: \texorpdfstring{$m = p$}{m=p} Prime \texorpdfstring{$\implies \mathbb{Z}/p\mathbb{Z}$}{=> Z/pZ} Is a Domain.}

Suppose $m = p$ is a prime number. Let $\bar{a}, \bar{b} \in \mathbb{Z}/p\mathbb{Z}$ and assume that their product vanishes:
\begin{equation}\label{eq:prime-prod-zero}
\bar{a} \cdot \bar{b} = \bar{0} \in \mathbb{Z}/p\mathbb{Z}.
\end{equation}
By definition of modular multiplication, this implies $\overline{a \cdot b} = \bar{0}$, meaning that $a \cdot b$ is congruent to $0 \pmod p$. Hence, there exists some integer $d \in \mathbb{Z}$ such that:
\begin{equation}\label{eq:divisible-by-p}
a \cdot b = d \cdot p \implies p \mid (a \cdot b).
\end{equation}
By Euclid's Lemma (which holds whenever $p$ is prime):
\begin{equation}\label{eq:euclids-lemma}
p \mid (a \cdot b) \implies p \mid a \quad \text{or} \quad p \mid b.
\end{equation}
Translating each case back into modular equivalence in $\mathbb{Z}/p\mathbb{Z}$:
\begin{itemize}
    \item If $p \mid a$, then $\bar{a} = \bar{0} \in \mathbb{Z}/p\mathbb{Z}$.
    \item If $p \mid b$, then $\bar{b} = \bar{0} \in \mathbb{Z}/p\mathbb{Z}$.
\end{itemize}
Thus, $\bar{a} \cdot \bar{b} = \bar{0}$ implies $\bar{a} = \bar{0}$ or $\bar{b} = \bar{0}$. This establishes that $\mathbb{Z}/p\mathbb{Z}$ has no zero divisors. Because $p \ge 2$, $\mathbb{Z}/p\mathbb{Z}$ contains at least two distinct elements ($\bar{0} \neq \bar{1}$), so $\mathbb{Z}/p\mathbb{Z}$ is not the zero ring. Therefore, $\mathbb{Z}/p\mathbb{Z}$ is an integral domain.
\end{content}

\begin{speech}[00:21:50 - 00:23:19]
So, you could say again in this proposition, when someone writes this, I here put $m$ to be a positive number ($m \in \mathbb{Z}_{>0}$), and now it matches carefully, but these boundary cases... you know, there's some aspect of this foundational part where you have to be careful yourself about boundary cases.

So, if I put $m = 1$, then I get the zero ring, but $1$ is not considered a prime number, so it matches. You could say: what happens if I put zero? It's not here; we don't think about primality of zero, so zero is left out of this, but if $m$ is zero, this is $\mathbb{Z}$, and it is a domain. Okay? All right.

\inlinemetanote{erases the blackboard}

Sometimes being too clear is actually more confusing in lectures. So, there's this kind of principle of lecturing where you lie about all the little details, and the lectures come out cleaner if you do that. And it's not clear didactically that that isn't a better way to do it.

There's a certain mathematician that takes pleasure in these edge cases, and you can tell that I'm not one of them.
\end{speech}
\end{proofWrapper}

\begin{content}[Pedagogical Analysis of Boundary Cases]
\begin{remark}[Boundary Cases: \texorpdfstring{$m = 1$}{m=1} and \texorpdfstring{$m = 0$}{m=0}]\label[remark]{rem:boundary-cases-mod-rings}
The proposition specifies $m \in \mathbb{Z}_{>0}$. Let us examine the boundary cases:
\begin{itemize}
    \item \textbf{Case $m = 1$:} The ideal is $1\mathbb{Z} = \mathbb{Z}$. The quotient ring is the zero ring:
    \[
    \mathbb{Z}/1\mathbb{Z} \cong \{0\},
    \]
    where $\bar{1} = \bar{0}$. Since the definition of an integral domain explicitly requires $R \neq \{0\}$, $\mathbb{Z}/1\mathbb{Z}$ is \emph{not} an integral domain. This perfectly matches the number-theoretic convention that $1$ is not a prime number.
    \item \textbf{Case $m = 0$:} The zero ideal is $0\mathbb{Z} = \{0\}$. The quotient ring is canonically isomorphic to the integers:
    \[
    \mathbb{Z}/0\mathbb{Z} \cong \mathbb{Z}.
    \]
    The integers $\mathbb{Z}$ are the prototypical integral domain. In modern commutative algebra, $\{0\} \subset \mathbb{Z}$ is a prime ideal precisely because $\mathbb{Z}/\{0\} \cong \mathbb{Z}$ is an integral domain!
\end{itemize}
\end{remark}
\end{content}

\section{Divisibility and Units in Commutative Rings}

\subsection{Divisibility in Arbitrary Rings}

\begin{speech}[00:23:19 - 00:24:29]
Okay, so something more fun is we get to talk about units, and this is actually pretty fun.

This has to do with every ring, so let's take a ring $R$.

And this is now related to last week; in fact, it was the point of the discussion last week. So, there we discussed the ring $\mathbb{Z}$, but now you can take any ring.

So, the definition here is that if we have elements $a$ and $b$ in the ring ($a, b \in R$), then we say --- this is the definition of the term ---, we say \qt{a divides b}. What this means is there exists some element of the ring, $c$ ($c \in R$), such that $a \cdot c = b$.

Okay, that's just the actual definition of division, and it's the same definition as with the integers.

So, this is a notion of division for any ring. Okay?
\end{speech}

\begin{content}[Divisibility Relations]
Let $R$ be a ring (assumed commutative with unity $1 \neq 0$).

\begin{definition}[Divisibility in a Ring]\label[definition]{def:divisibility-ring}
Let $a, b \in R$. We say that \newterm{$a$ divides $b$}, denoted:
\begin{equation}\label{eq:divisibility-notation}
a \mid b,
\end{equation}
if there exists an element $c \in R$ such that:
\begin{equation}\label{eq:divisibility-def}
a \cdot c = b.
\end{equation}
In this case, we also say that $b$ is a \newterm{multiple} of $a$, or that $a$ is a \newterm{divisor} of $b$.
\end{definition}

\begin{steps}[Connection with Principal Ideals]
In ideal-theoretic language, $a \mid b$ in a commutative ring $R$ is equivalent to the statement that the principal ideal generated by $b$ is contained in the principal ideal generated by $a$:
\[
a \mid b \iff b \in aR \iff (b) \subseteq (a).
\]
\end{steps}
\end{content}

\subsection{The Definition of a Unit}

\begin{speech}[00:24:29 - 00:26:22]
Now, $u$ in $R$ ($u \in R$) is a unit --- that's the definition of a unit.

The definition is if $u$ divides $1$ ($u \mid 1$). Oh, I should say that also: \qt{a divides b}, that's how you say it in English, and the symbol for that is $a \mid b$. Okay?

So, here this is now a definition: an element is a unit if $u$ divides $1$ ($u \mid 1$). And that is equivalent --- this is just equivalent --- to there existing a $w$ ($w \in R$) such that $u \cdot w = 1$. I never have to worry about the order here because I said that for the rings in this course, unless I tell you explicitly otherwise, we're always in commutative rings. Okay?

So, this is the definition, this is unwinding division, and if I want to unwind this using English --- so this is the definition of division, and this is just English ---, we say that if and only if $u$ has a multiplicative inverse in $R$.

Okay? So, this is a really fundamental definition: a unit in our ring. So, are you happy with this? That's the definition, this is just translation, and that's what you say: $u$ has a multiplicative inverse, $w$ is the multiplicative inverse. So, is there any problem with this?
\end{speech}

\begin{content}[Invertible Elements]
\begin{definition}[Unit in a Ring]\label[definition]{def:unit-ring}
Let $R$ be a commutative ring with unity $1$. An element $u \in R$ is called a \newterm{unit} if $u$ divides $1$:
\begin{equation}\label{eq:unit-div-1}
u \mid 1.
\end{equation}
By the definition of divisibility (\cref{def:divisibility-ring}), this is equivalent to:
\begin{equation}\label{eq:unit-mult-inv}
\exists w \in R \quad \text{such that} \quad u \cdot w = 1.
\end{equation}
In words, $u$ is a unit if and only if $u$ has a multiplicative inverse in $R$. The element $w$ is uniquely determined and denoted $u^{-1} := w$.
\end{definition}

\begin{steps}[Equivalence of Characterizations on the Blackboard]
The blackboard outlines three equivalent perspectives on units:
\begin{enumerate}[label=\textbf{(\arabic*)}]
    \item \textbf{Divisibility Perspective:} $u \mid 1$ (the element divides the multiplicative identity).
    \item \textbf{Algebraic Perspective (Unwinding Division):} $\exists w \in R$ such that $u \cdot w = 1$.
    \item \textbf{Verbal Perspective (English):} $u$ possesses a multiplicative inverse in $R$.
\end{enumerate}
Since all rings considered here are commutative, left-inverses and right-inverses coincide.
\end{steps}
\end{content}

\begin{hidden}
% timestamp: 00:26:22
% topic: Definition of divisibility and units in arbitrary commutative rings
% board_state: def:divisibility-ring, def:unit-ring
% next_goal: Linguistic distinction between "the unit", "the units", and "a unit", followed by the units of Z
% open_loops: none
\end{hidden}

\subsection{Linguistic Nuances: \qt{The Units} vs. \qt{The Unit} vs. \qt{A Unit}}

\begin{speech}[00:26:22 - 00:28:16]
Actually, people asked me some questions about language, and you know, English is a language I know, which means that I speak it without thinking. It was pointed out to me that the use of \qt{unit} is subtle, and I was doing it already in the first week; if you were sensitive to language, you'd have seen. So, for $\mathbb{Z}$ here, the units: what are the units for $\mathbb{Z}$?

In $\mathbb{Z}$, and it comes as no surprise, this is $+1$ and $-1$. Those are the only invertibles in it.

So, the units: if I say \qt{the units} here, that means all of them. The units.

If I say \qt{the unit} --- this is standard use of mathematical English ---, if I say \qt{the unit}, that's just $1$. If I have a ring and say \qt{let's consider the unit}, that means $1$.

You could say \qt{a unit}. If I say \qt{let's consider a ring, and let's consider a unit}, that could be either one; we don't know. Could be either. A unit --- this could be either.

So, that's actually standard usage. It's a little confusing, but once you get used to it, you know, it's just language. Is it clear? It was pointed out to me that I was doing that, and anyway, that's right. But if someone comes to you with a ring and says \qt{consider the unit}, that person means $1$.
\end{speech}

\begin{content}[Units of the Integers and Linguistic Distinctions]
\begin{example}[Units of the Ring of Integers \texorpdfstring{$\mathbb{Z}$}{Z}]\label[example]{ex:units-of-integers}
In the ring of integers $\mathbb{Z}$, the only elements possessing multiplicative inverses are $1$ and $-1$:
\begin{equation}\label{eq:units-of-integers}
\text{The units of } \mathbb{Z} = \{1, -1\}.
\end{equation}
\end{example}

\begin{remark}[Linguistic Distinction in Mathematical English]\label[remark]{rem:unit-terminology}
In mathematical English, the word \qt{unit} carries three distinct grammatical usages depending on the article:
\begin{itemize}
    \item \textbf{\qt{The units} (Plural):} Refers to the entire set of all invertible elements in the ring $R$, denoted $U(R)$ or $R^\times$. For $\mathbb{Z}$, the units are $\{1, -1\}$.
    \item \textbf{\qt{The unit} (Singular, Definite Article):} Refers specifically to the multiplicative identity element $1_R \in R$.
    \item \textbf{\qt{A unit} (Singular, Indefinite Article):} Refers to an arbitrary invertible element $u \in R$ (i.e., any element $u \in U(R)$). For $\mathbb{Z}$, \qt{a unit} could be either $1$ or $-1$.
\end{itemize}
\end{remark}
\end{content}

\section{The Group of Units and Modular Arithmetic}

\subsection{The Group of Units of a Ring}

\begin{speech}[00:28:16 - 00:29:11]
Okay, so we discussed the units of the Gaussian integers (Recall: $\pm 1, \pm i$, Week~1, \cref{def:gaussian-integers}), and now, since everything is defined, I recommend proving what the units of the Gaussian integers are. We wrote them last time.

But let's do something else now.

So, this is an exercise... \inlinemetanote{writes on board} Exercise: prove the units of the Gaussian integers are $\pm 1, \pm i$. You need an argument for that, and that argument comes from the norm in complex analysis, so prove it.

Okay, so observation...
\end{speech}

\begin{content}[Exercise: Units of the Gaussian Integers]
\begin{exercise}[Units of the Gaussian Integers]\label[exercise]{exer:units-gaussian-integers}
Prove that the units of the ring of Gaussian integers $\mathbb{Z}[i] = \{a + bi \mid a, b \in \mathbb{Z}\}$ are precisely:
\begin{equation}\label{eq:units-gaussian-integers}
U(\mathbb{Z}[i]) = \{\pm 1, \pm i\}.
\end{equation}
\end{exercise}

The lecturer writes out this proof in Week~3 (\cref{prop:units-gaussian-integers}).

\begin{steps}[Proof Strategy via the Multiplicative Norm]
Consider the field norm $N \colon \mathbb{Q}(i) \to \mathbb{Q}$ defined by $N(a + bi) = a^2 + b^2$. For all $z, w \in \mathbb{Z}[i]$, the norm is multiplicative:
\[
N(z \cdot w) = N(z) \cdot N(w).
\]
If $z \in \mathbb{Z}[i]$ is a unit, there exists $w \in \mathbb{Z}[i]$ with $z \cdot w = 1$. Taking norms yields:
\[
N(z) \cdot N(w) = N(1) = 1.
\]
Since $a, b \in \mathbb{Z}$, $N(z) = a^2 + b^2$ is a non-negative integer. The only positive integer divisor of $1$ in $\mathbb{Z}$ is $1$, so $a^2 + b^2 = 1$. The only integer solutions are:
\[
(a, b) \in \{(\pm 1, 0), (0, \pm 1)\} \implies z \in \{\pm 1, \pm i\}.
\]
\end{steps}
\end{content}

\begin{speech}[00:29:11 - 00:30:47]
...and it's a really important observation: if I take any ring $R$ --- let $R$ be a ring --- and let $U$ --- you could write $U(R)$ if you want --- be the set of units of $R$.

Then this observation is: if I take the units, I take multiplication --- that's multiplication in the ring ---, and I take the unit $1$, this is a group.

Okay? What does it say? If I take two units and multiply them, I get a unit.

And that's clear because, you know, if you have two elements here, if I multiply them, I have to show that's a unit. That means I have to show that if I multiply by some more stuff, I get $1$.

But this is a unit, so there's already its multiplicative inverse, and that one is a unit, so there's already its multiplicative inverse; so if I multiply these together, I get $1$, because I can flip the orders around.

So, this says the units are closed under multiplication.

And there's the neutral element, and then you have to show the existence of inverses, but a unit is a unit because it has an inverse.

So, that's the proof of this. This is a good thing to think about: this is always a group.
\end{speech}

\begin{content}[The Group of Units]
\begin{proposition}[Group of Units]\label[proposition]{prop:group-of-units}
Let $R$ be a ring with unity $1_R$, and let $U(R)$ (also denoted $U_R$ or $R^\times$) be the set of units of $R$:
\begin{equation}\label{eq:unit-set-def}
U(R) := \{u \in R \mid \exists w \in R \colon u \cdot w = 1_R\}.
\end{equation}
Then $(U(R), \cdot, 1_R)$ is a group under the multiplication of $R$.
\end{proposition}

\begin{proof}
We verify the group axioms for $(U(R), \cdot, 1_R)$:
\begin{itemize}
    \item \textbf{Closure under Multiplication:} Let $u_1, u_2 \in U(R)$. By definition, there exist $w_1, w_2 \in R$ such that $u_1 w_1 = 1_R$ and $u_2 w_2 = 1_R$. Since $R$ is commutative:
    \[
    (u_1 u_2)(w_2 w_1) = (u_1 w_1)(u_2 w_2) = 1_R \cdot 1_R = 1_R.
    \]
    Thus, $u_1 u_2$ has a two-sided multiplicative inverse in $R$, which shows $u_1 u_2 \in U(R)$.
    \item \textbf{Associativity:} Multiplicative associativity in $U(R)$ is directly inherited from the ring $R$.
    \item \textbf{Identity Element:} Since $1_R \cdot 1_R = 1_R$, the unity $1_R$ is its own inverse, so $1_R \in U(R)$. For all $u \in U(R)$, $1_R \cdot u = u$.
    \item \textbf{Existence of Inverses:} For any $u \in U(R)$, there exists $w \in R$ such that $u w = 1_R$. By symmetry of the relation $w u = 1_R$, $w$ is also invertible with inverse $u$, so $w = u^{-1} \in U(R)$.
\end{itemize}
Hence, $(U(R), \cdot, 1_R)$ is a group in the sense of \cref{def:group}. If $R$ is commutative, $U(R)$ is an abelian group.
\end{proof}
\end{content}

\subsection{Characterization of Units in \texorpdfstring{$\mathbb{Z}/m\mathbb{Z}$}{Z/mZ}}

\begin{speech}[00:30:47 - 00:31:44]
So, a very natural question, and one of the first things to think about regarding $\mathbb{Z}/m\mathbb{Z}$, if I take this ring $R$ \inlinemetanote{writes on board}...

And it was implicit there: I will always consider $m > 0$ when I do this, because there's no reason to consider negative ones, as that's just counting everything twice. And then there's the confusing case when $m = 0$, which is usually a special case, so I exclude it.

Usually we allow $m = 1$, which is the zero ring, but anyway, this is kind of standard. When I write $\mathbb{Z}/m\mathbb{Z}$, I always mean this, and typically I also mean $m > 1$, but anyway.

So, the question is: what are the units of $\mathbb{Z}/m\mathbb{Z}$?
\end{speech}

\begin{speech}[00:31:44 - 00:33:33]
So, let's say: suppose we have $\bar{x} \in \mathbb{Z}/m\mathbb{Z}$, and $\bar{x}$ is a unit.

This means that there exists $\bar{y} \in \mathbb{Z}/m\mathbb{Z}$ such that $\bar{x} \cdot \bar{y} = \bar{1}$ in $\mathbb{Z}/m\mathbb{Z}$.

I can write this, or sometimes you write $\bmod 1$, but anyway, that's what it means to be a unit --- I'm just repeating it, right?

And then I can translate: what does this mean in this ring? What does it mean for the integers $x$, $y$, and $1$? This is equivalent to $x \cdot y = 1 + d \cdot m$.

Do you agree with that? That's what it means, for some $d$. Okay?

And this, if you remember from the first week, means that $x$ and $m$ have $\operatorname{gcd} = 1$.

The only thing that could divide both $x$ and $m$, it would also have to divide $1$. So, that's why we say $\operatorname{gcd} = 1$.

In English, this condition is usually said as $x$ and $m$ are relatively prime.

Those are just words that express $\operatorname{gcd}(x, m) = 1$. Okay?
\end{speech}

\begin{hidden}
% timestamp: 00:33:33
% topic: Units of Z/mZ characterized via coprimality (gcd(x, m) = 1)
% board_state: prop:group-of-units, thm:units-z-mod-m
% next_goal: Formulate explicit theorem that U(Z/mZ) consists of coprime residue classes and is a group
% open_loops: none
\end{hidden}

\begin{speech}[00:33:33 - 00:35:02]
So, we have now discovered what the units of $\mathbb{Z}/m\mathbb{Z}$ are: they are all the elements which are relatively prime.

Because if it is a unit, it's the residue class of something relatively prime; if you have $\operatorname{gcd} = 1$, then you go backwards in this diagram.

So, I state that as a result.

So, if you write this as a result, then it says that the units of $\mathbb{Z}/m\mathbb{Z}$ consist of all elements $\bar{x}$ where $\operatorname{gcd}(x, m) = 1$. Did I write $\operatorname{gcd}$ here? Maybe $\operatorname{gcd}$. Okay? Is that clear?

And so the thing we've learned is that this units group, $U(\mathbb{Z}/m\mathbb{Z})$, is a group.

So, it says in particular, if you take $\mathbb{Z}/m\mathbb{Z}$, and you take the residue classes that are relatively prime, they form a group under multiplication.

So, let me do an example.
\end{speech}

\begin{content}[Units of the Quotient Ring \texorpdfstring{$\mathbb{Z}/m\mathbb{Z}$}{Z/mZ}]
\begin{theorem}[Units in $\mathbb{Z}/m\mathbb{Z}$]\label[theorem]{thm:units-z-mod-m}
Let $m \in \mathbb{Z}_{>0}$ be a positive integer. An element $\bar{x} \in \mathbb{Z}/m\mathbb{Z}$ is a unit if and only if $x$ and $m$ are coprime (this does not depend on the representative $x$, since $\operatorname{gcd}(x + km, m) = \operatorname{gcd}(x, m)$):
\begin{equation}\label{eq:units-z-mod-m}
U(\mathbb{Z}/m\mathbb{Z}) = (\mathbb{Z}/m\mathbb{Z})^\times = \left\{ \bar{x} \in \mathbb{Z}/m\mathbb{Z} \;\big|\; \operatorname{gcd}(x, m) = 1 \right\}.
\end{equation}
The order of this group is given by Euler's totient function $\phi(m) := |U(\mathbb{Z}/m\mathbb{Z})|$ (named by the lecturer in Week~3, \cref{def:eulers-totient}).
\end{theorem}

\begin{proof}
We establish the equivalence by unwinding the definition of divisibility:
\begin{itemize}
    \item \textbf{Forward Direction ($\implies$):} Suppose $\bar{x} \in U(\mathbb{Z}/m\mathbb{Z})$. Then there exists $\bar{y} \in \mathbb{Z}/m\mathbb{Z}$ such that:
    \[
    \bar{x} \cdot \bar{y} = \bar{1} \in \mathbb{Z}/m\mathbb{Z}.
    \]
    In terms of integer divisibility, this means $x y \equiv 1 \pmod m$. Hence, there exists an integer $d \in \mathbb{Z}$ such that:
    \[
    x y = 1 + d \cdot m \iff x y - d \cdot m = 1.
    \]
    Any common divisor $g \in \mathbb{Z}$ of $x$ and $m$ must divide the linear combination $x y - d m = 1$. Therefore, $g \mid 1$, which forces $\operatorname{gcd}(x, m) = 1$.
    \item \textbf{Reverse Direction ($\impliedby$):} Conversely, suppose $\operatorname{gcd}(x, m) = 1$. By Bézout's Identity, there exist integers $y, k \in \mathbb{Z}$ such that:
    \[
    x \cdot y + m \cdot k = 1 \implies x \cdot y \equiv 1 \pmod m.
    \]
    Projecting this relation into residue classes in $\mathbb{Z}/m\mathbb{Z}$, we obtain $\bar{x} \cdot \bar{y} = \bar{1}$. Thus, $\bar{x}$ is invertible, so $\bar{x} \in U(\mathbb{Z}/m\mathbb{Z})$.
\end{itemize}
By \cref{prop:group-of-units}, $(U(\mathbb{Z}/m\mathbb{Z}), \cdot, \bar{1})$ forms a finite abelian group.
\end{proof}
\end{content}

\subsection{Example: The Unit Group of \texorpdfstring{$\mathbb{Z}/7\mathbb{Z}$}{Z/7Z}}

\begin{speech}[00:35:02 - 00:36:09]
$\mathbb{Z}/7\mathbb{Z}$.

So, that's actually a prime one, but let me now do this example. So, the elements of this are $\bar{0}, \bar{1}, \bar{2}, \bar{3}, \bar{4}, \bar{5}, \bar{6}$, right?

And which of these residue classes come here? The classes of $x$ where $x$ and $7$ are relatively prime.

So, the question is: what are the units? Of course, it's some of these, but which of these residue classes correspond to $x$ and $7$ being relatively prime?

What would you say? As I said, if you don't know, guess.
\end{speech}

\begin{student}[Student Answer]
All of them.
\end{student}

\begin{speech}[00:36:09 - 00:36:14]
All of them? What about this one \inlinemetanote{points at $\bar{0}$}?
\end{speech}

\begin{student}[Student Answer]
Only $1$. $1$ to $6$ are the units.
\end{student}

\begin{speech}[00:36:14 - 00:36:50]
Exactly. So, the answer to this question is these: $\{\bar{1}, \bar{2}, \bar{3}, \bar{4}, \bar{5}, \bar{6}\}$.

And the reason is because, look, this was the answer: it's the residue classes of those numbers that are relatively prime. So, this is the residue class of those numbers that are relatively prime to $7$. And these are those residue classes.

If you're $0 \bmod 7$, then you're not relatively prime to $7$. Like an example of $0 \bmod 7$ would be $14$, and that's not relatively prime to $7$. So, it's these. Is that clear?
\end{speech}

\begin{content}[The Unit Group of \texorpdfstring{$\mathbb{Z}/7\mathbb{Z}$}{Z/7Z}]
\begin{example}[Units of $\mathbb{Z}/7\mathbb{Z}$]\label[example]{ex:units-z7z}
Consider the quotient ring $\mathbb{Z}/7\mathbb{Z}$. The elements of the ambient ring are:
\[
\mathbb{Z}/7\mathbb{Z} = \{\bar{0}, \bar{1}, \bar{2}, \bar{3}, \bar{4}, \bar{5}, \bar{6}\}.
\]
Since $7$ is a prime number, every non-zero residue class $\bar{x} \in \{\bar{1}, \bar{2}, \bar{3}, \bar{4}, \bar{5}, \bar{6}\}$ satisfies $\operatorname{gcd}(x, 7) = 1$. The zero class $\bar{0}$ is not coprime to $7$ (since $\operatorname{gcd}(0, 7) = 7 \neq 1$, or for example $14 \equiv 0 \pmod 7$ where $\operatorname{gcd}(14, 7) = 7$).
Consequently, the group of units is:
\begin{equation}\label{eq:units-z7z-set}
U(\mathbb{Z}/7\mathbb{Z}) = (\mathbb{Z}/7\mathbb{Z})^\times = \{\bar{1}, \bar{2}, \bar{3}, \bar{4}, \bar{5}, \bar{6}\}.
\end{equation}
The order of the unit group is $|U(\mathbb{Z}/7\mathbb{Z})| = \phi(7) = 7 - 1 = 6$.
\end{example}
\end{content}

\subsubsection{Powers of the Element \texorpdfstring{$a = \bar{3}$}{a = 3} and Cyclic Structure}

\begin{speech}[00:36:50 - 00:37:09]
Okay, and the claim is this is a group under multiplication. Well, anyway, this is a group under addition, but if you take the relatively prime ones, this is a group under multiplication.

So, we should be able to figure out what group it is. So, what's the neutral element of this group?
\end{speech}

\begin{student}[Student Answer]
One.
\end{student}

\begin{speech}[00:37:09 - 00:37:56]
One, it's going to be $\bar{1}$.

What is this group? It's a commutative group of order $6$. And so later we'll see that we know all commutative groups of order $6$. We'd like to know which one this is. It's isomorphic to some group, we know.

And so it matters what you pick. I propose we pick $3$. So, let me say we take this element in this group, $a = \bar{3}$.

And let's try to consider its powers. So, there's the neutral element, then I can take $a, a^2, a^3, a^4$. You should know the powers of $3$. Okay, let's see how many powers of $3$ you know.
\end{speech}

\begin{speech}[00:37:56 - 00:39:44]
So, what is this one? So I'm just going to do this: this is $\bar{1}$. So, what is $a$? Well, that's just that: $\bar{3}$.

What is this? That's the square of $a$, so that's the residue class of $3^2$, which is $9$. So, I should put $\bar{2}$ here, right? Are you happy with this?

And this is $a^3$: $3^3$ is $27$. For $27$, the residue $\bmod 7$ is $\bar{6}$, right?

$a^4$ is $3^4$, which is $81$, and the residue $\bmod 7$ of $81$ is $\bar{4}$, I think.

$3^5$ is $243$, and its residue is going to be $\bar{5}$.

What do you think this is going to be? Do you know what $3^6$ is? $3^6$ is $729$. And if I look at its residue class $\bmod 7$, well, $700$ is already here, and $28$ is divisible, so its residue class is $\bar{1}$. This is back to $1$.

So, this is understanding the group. You should know the powers of $3$; they're very useful, you know.

So, this is what happens when you multiply in this group: you take $a$, $a^2$, the next one $a^3$, the next one $a^4$, the next one $a^5$, and then when you go to the sixth, you go back to the identity.

What group do you know this happens in? Does this remind you of anything? Yeah?
\end{speech}

\begin{student}[Student Answer]
A cyclic group.
\end{student}

\begin{speech}[00:39:44 - 00:39:54]
It's $\mathbb{Z}/6\mathbb{Z}$. So, the proposition that we've kind of proven here...
\end{speech}

\begin{content}[Powers of the Generator \texorpdfstring{$a = \bar{3}$}{a = 3} in \texorpdfstring{$U(\mathbb{Z}/7\mathbb{Z})$}{U(Z/7Z)}]
To examine the internal structure of the abelian group $U(\mathbb{Z}/7\mathbb{Z})$ of order $6$, we choose the candidate generator $a := \bar{3} \in U(\mathbb{Z}/7\mathbb{Z})$ and compute its successive powers:

\begin{center}
\begin{tabular}{c|ccccccc}
\toprule
\textbf{Power} & $a^0$ & $a^1$ & $a^2$ & $a^3$ & $a^4$ & $a^5$ & $a^6$ \\
\midrule
\textbf{Symbolic} & $\bar{1}$ & $a$ & $a^2$ & $a^3$ & $a^4$ & $a^5$ & $a^6$ \\
\textbf{Value in $\mathbb{Z}$} & $1$ & $3$ & $9$ & $27$ & $81$ & $243$ & $729$ \\
\textbf{Residue modulo $7$} & $\bar{1}$ & $\bar{3}$ & $\bar{2}$ & $\bar{6}$ & $\bar{4}$ & $\bar{5}$ & $\bar{1}$ \\
\bottomrule
\end{tabular}
\end{center}

\begin{steps}[Modular Arithmetic Derivations for Powers of $3$]
\begin{align*}
    a^0 &= \bar{1}, \\
    a^1 &= \bar{3}, \\
    a^2 &= \overline{3^2} = \bar{9} = \bar{2} && (9 = 7 \cdot 1 + 2), \\
    a^3 &= \overline{3^3} = \overline{27} = \bar{6} && (27 = 7 \cdot 3 + 6), \\
    a^4 &= \overline{3^4} = \overline{81} = \bar{4} && (81 = 7 \cdot 11 + 4), \\
    a^5 &= \overline{3^5} = \overline{243} = \bar{5} && (243 = 7 \cdot 34 + 5), \\
    a^6 &= \overline{3^6} = \overline{729} = \bar{1} && (729 = 7 \cdot 104 + 1).
\end{align*}
Because the powers of $a = \bar{3}$ exhaust all $6$ elements of $U(\mathbb{Z}/7\mathbb{Z})$ before returning to the identity $\bar{1}$ at exponent $6$, the element $\bar{3}$ is a \newterm{generator} (or \newterm{primitive root} modulo $7$). Consequently, $U(\mathbb{Z}/7\mathbb{Z})$ is a cyclic group of order $6$.
\end{steps}
\end{content}

\begin{hidden}
% timestamp: 00:39:54
% topic: Cyclic generator of U(Z/7Z) and establishment of powers of 3 mod 7
% board_state: ex:units-z7z, generator-powers-table
% next_goal: Formalize the group isomorphism U(Z/7Z) \cong (Z/6Z, +, 0) via explicit discrete logarithm map
% open_loops: none
\end{hidden}

\subsubsection{The Group Isomorphism \texorpdfstring{$(U(\mathbb{Z}/7\mathbb{Z}), \cdot) \cong (\mathbb{Z}/6\mathbb{Z}, +)$}{(U(Z/7Z), *) = (Z/6Z, +)}}

\begin{speech}[00:39:54 - 00:41:24]
...we've proven what that calculation of powers of $3$... and I picked $3$. If you picked another one, it would have been slightly different, but that calculation proves a group isomorphism, which maybe I'll write here and we think about.

You should check my powers of $3$, and then after you take the power of $3$, you have to take the residue $\bmod 7$.

So, this great 20th century Hungarian mathematician (Paul Erd\H{o}s) said that if numbers aren't beautiful, I don't know what is. That's kind of the mathematical perspective on numbers.

So, we have here the units of $\mathbb{Z}/7\mathbb{Z}$, and then there's another group, which is $\mathbb{Z}/6\mathbb{Z}$ with $+$ and $\bar{0}$.

And I claim these are isomorphic. So, this is the units with multiplication and $\bar{1}$. This group is isomorphic to that group. And what is the isomorphism? It's written on the board.

The isomorphism --- this is a slightly painful way to write the isomorphism, but let's just be painful about it ---: so, here $\bar{1}$ is going to go to $\bar{0}$. And bar means something different: it means $\bmod 7$ here, and $\bmod 6$ here, right?
\end{speech}

\begin{speech}[00:41:24 - 00:42:33]
$a$, which is $3$ here, $a$ is $\bar{3}$. So, the nice way to do it is: $a$ goes to $\bar{1}$, $a^2$ goes to $\bar{2}$. I guess I could write this in a slightly better notation, but anyway, you get the pattern. There are not so many more: $a^3, a^4, a^5$. These go to $\bar{3}, \bar{4}, \bar{5}$.

So, this is the simple way to write it. And the reason this is true is because on this side when we multiply, we just add up the exponents, and on this side we add up those. And the important thing is the sixth one goes to zero, which is what we have here, okay?

And if you want to say this explicitly, you have to write these $a$'s in terms of the elements of the units, which I've done there. So, this is totally explicit.

Okay, that's all. \inlinemetanote{students applaud and pack up; video ends}
\end{speech}

\begin{content}[Group Isomorphism: \texorpdfstring{$U(\mathbb{Z}/7\mathbb{Z}) \cong (\mathbb{Z}/6\mathbb{Z}, +, \bar{0})$}{U(Z/7Z) isomorphic to (Z/6Z, +, 0)}]
\begin{proposition}[Cyclic Group Isomorphism]\label[proposition]{prop:z7z-cyclic-isomorphism}
The multiplicative group of units $(U(\mathbb{Z}/7\mathbb{Z}), \cdot, \bar{1})$ is isomorphic to the additive cyclic group $(\mathbb{Z}/6\mathbb{Z}, +, \bar{0})$:
\begin{equation}\label{eq:z7z-isom-group}
(U(\mathbb{Z}/7\mathbb{Z}), \cdot, \bar{1}) \cong (\mathbb{Z}/6\mathbb{Z}, +, \bar{0}).
\end{equation}
\end{proposition}

\begin{proof}
Let $a = \bar{3} \in U(\mathbb{Z}/7\mathbb{Z})$. We define the mapping $\Phi \colon U(\mathbb{Z}/7\mathbb{Z}) \to \mathbb{Z}/6\mathbb{Z}$ by sending powers of the generator $a$ to their integer exponents modulo $6$:
\begin{align*}
    \Phi \colon U(\mathbb{Z}/7\mathbb{Z}) &\to \mathbb{Z}/6\mathbb{Z}, \\
    \bar{1} = a^0 &\mapsto \bar{0}, \\
    \bar{3} = a^1 &\mapsto \bar{1}, \\
    \bar{2} = a^2 &\mapsto \bar{2}, \\
    \bar{6} = a^3 &\mapsto \bar{3}, \\
    \bar{4} = a^4 &\mapsto \bar{4}, \\
    \bar{5} = a^5 &\mapsto \bar{5}.
\end{align*}
\textbf{Homomorphism Property:} For any exponents $j, k \in \{0, 1, 2, 3, 4, 5\}$, multiplication in the unit group corresponds to addition of exponents:
\[
a^j \cdot a^k = a^{j+k} = a^{(j+k) \bmod 6},
\]
since $a^6 = \bar{1}$. Applying the map $\Phi$:
\[
\Phi(a^j \cdot a^k) = \overline{j + k} = \bar{j} + \bar{k} = \Phi(a^j) + \Phi(a^k) \in \mathbb{Z}/6\mathbb{Z}.
\]
Hence, $\Phi$ is a group homomorphism. By the table of powers, $a^0, \dots, a^5$ are exactly the six units, and $\Phi$ sends them to the six distinct classes $\bar{0}, \dots, \bar{5}$; so $\Phi$ is bijective, which establishes the isomorphism.
\end{proof}

\begin{steps}[Disambiguating Residue Class Notations Across Groups]
The bar notation $\bar{\cdot}$ operates over two distinct moduli in this isomorphism:
\begin{itemize}
    \item In the domain $(U(\mathbb{Z}/7\mathbb{Z}), \cdot)$, bars denote residue classes modulo $7$ under multiplication: $\bar{x} = x + 7\mathbb{Z}$.
    \item In the codomain $(\mathbb{Z}/6\mathbb{Z}, +)$, bars denote residue classes modulo $6$ under addition: $\bar{k} = k + 6\mathbb{Z}$.
\end{itemize}
The isomorphism $\Phi$ is the \newterm{discrete logarithm} with respect to the base $3$ modulo $7$.
\end{steps}
\end{content}

\begin{insight}[Paul Erd\H{o}s on the Beauty of Numbers]
The lecturer references the celebrated quote by the Hungarian mathematician Paul Erd\H{o}s (1913--1996): \qt{If numbers are not beautiful, I don't know what is.} The unexpected isomorphism between the multiplicative units of a prime field $\mathbb{F}_7^\times$ and the simple additive clock arithmetic of $\mathbb{Z}/6\mathbb{Z}$ is a classic embodiment of this aesthetic: multiplication modulo $7$ transforms into addition modulo $6$ via the powers of a single primitive root.
\end{insight}

\begin{overview}[Summary of Key Concepts]
\begin{itemize}
    \item \textbf{Category Theory:} A category $\mathcal{C}$ comprises objects $\operatorname{Obj}(\mathcal{C})$, hom-sets $\operatorname{Hom}(A, B)$, associative composition $\circ$, and identity morphisms $\operatorname{Id}_A$.
    \item \textbf{Categorical Morphisms:}
    \begin{itemize}
        \item \emph{Isomorphism:} Invertible morphism ($f \circ g = \operatorname{Id}_B$, $g \circ f = \operatorname{Id}_A$).
        \item \emph{Monomorphism:} Left-cancellable morphism ($f \circ g_1 = f \circ g_2 \implies g_1 = g_2$), generalizing injectivity.
        \item \emph{Epimorphism:} Right-cancellable morphism ($g_1 \circ f = g_2 \circ f \implies g_1 = g_2$), generalizing surjectivity.
    \end{itemize}
    \item \textbf{Integral Domains:} A commutative ring with $1 \neq 0$ having no zero divisors ($a \neq 0 \land b \neq 0 \implies ab \neq 0$). The quotient ring $\mathbb{Z}/m\mathbb{Z}$ is an integral domain if and only if $m$ is prime.
    \item \textbf{Units and Divisibility:} An element $u \in R$ is a unit if $u \mid 1$, i.e., $u$ has a multiplicative inverse. The units form a multiplicative group $U(R)$.
    \item \textbf{Units of Residue Class Rings:} $U(\mathbb{Z}/m\mathbb{Z}) = \{ \bar{x} \in \mathbb{Z}/m\mathbb{Z} \mid \operatorname{gcd}(x, m) = 1 \}$ with order $\phi(m)$. For prime $p = 7$, $U(\mathbb{Z}/7\mathbb{Z}) \cong (\mathbb{Z}/6\mathbb{Z}, +)$ generated by $\bar{3}$.
\end{itemize}
\end{overview}